{\large\textbf{Photonic Crystal (20 points, 5 hours)}}

\textbf{Introduction}

Photonic crystals are materials with periodic variation of refractive index on
a scale comparable to the wavelength of light. In the optical spectrum of
photonic crystals, narrow regions of wavelengths exist within which the light
propagation is suppressed. These unusual optical properties are used for the
development of diverse optical elements (optical filters, reflectors) on the
basis of photonic crystals.

In this research work, you are invited to study the properties of photonic
crystals based on porous (containing holes filled with air) films of anodic
aluminum oxide. The structure of porous films of anodic aluminum oxide prepared
by electrochemical oxidation (anodization) of aluminum can be represented as a
system of disjoint cylindrical channels which are located perpendicular to the
film surface (Fig. 1).

\begin{center}\includegraphics[width=0.9\linewidth]{figures/APhO_2017_Q4_fig1.jpg}\end{center}

Fig 1: (a) Schematic structure of the porous film of anodic aluminum oxide: (1)
pore, (2) porous layer, (3) barrier layer. Electron microscopy image of the
film: (b) top view, (c) crosssection [Electrochim. Acta, 2011, 56, 2378].

Pore diameter and interpore distance depend on the conditions of
electrochemical treatment allowing one to obtain structures with variable
porosity along the normal to the surface of the oxide film by changing the
voltage during anodization (Fig. 2).

\begin{center}\includegraphics[width=0.36\linewidth]{figures/APhO_2017_Q4_fig2.jpg}\end{center}

Fig 2: Electron microscopy image of cross section of anodic alumina with
variable pore diameter [Nat.Mater., 2006, 5, 741].

The pore diameter of the samples studied in this work is less than 30\,nm.
Porous medium with such a small pore diameter becomes continuous for the
electromagnetic waves within the optical range and, as a consequence, it can be
characterized by effective (volume-average) refraction index.

Variation of the oxide film porosity leads to the variation of the refractive
index. It should be noted that the refractive index of the studied samples
periodically changes only in a single direction: along the normal to the
surface of the film. Therefore, the studied samples are \textit{one-dimensional
photonic crystals}.

It should be pointed out that the thin films of anodic aluminum oxide with
uniform diameter are optically transparent whereas photonic crystals become
colored due to the light interference within their layered structure.

During this research, you will perform optical studies of three samples of
photonic crystals based on anodic aluminum oxide (AAO) with the structures of
various complexities.

\textbf{Equipment}

\begin{center}\includegraphics[width=0.9\linewidth]{figures/APhO_2017_Q4_fig3.jpg}\end{center}

\begin{center}Fig. 3. Equipment on the table.\end{center}

1. Electrical outlet

2. Diffraction arm for spectral measurements

3. Laser holder with rubber band

4. Photodiode with holder

5. Sample holder

6. Main optical bench with protractors

7. Water in bottle

8. DC adapter (power supply)

9. Red laser (wavelength 659\,nm)

10. Green laser (wavelength 530\,nm)

11. Blue laser (wavelength 400\,nm)

\noindent\fbox{\parbox{\dimexpr\linewidth-2\fboxsep-2\fboxrule\relax}{\centering\large
Lasers are dangerous! Do not point lasers at your eyes or look into the laser
beam under any circumstances! Beware the reflected beams!}}

12. Laser button lock

13. Petri dish with three photonic crystal samples in square color frames

\noindent\fbox{\parbox{\dimexpr\linewidth-2\fboxsep-2\fboxrule\relax}{\centering\large
Photonic crystals are very fragile! Do not touch the crystals! Hold the samples
only by frame! New samples will not be provided!}}

14. Petri dish with the coverslip

\noindent\fbox{\parbox{\dimexpr\linewidth-2\fboxsep-2\fboxrule\relax}{\centering\large
Glass is thin and fragile! Use tweezers to hold it!}}

15. Tweezers to operate the coverslip

16. Syringe

17. Multimeter with connector wires

18. Ruler

\textbf{Diffraction arm}

For spectral measurements, you will need diffraction arm (fig. 4) to obtain
different wavelengths. Diffraction arm is composed of the filament lamp (L), a
diffraction grating (DG) in the holder and a movable lens (F) between them. To
install the diffraction arm place the disk into the left hole of the main
bench.

\begin{center}\includegraphics[width=0.6\linewidth]{figures/APhO_2017_Q4_fig4.jpg}\end{center}

\begin{center}Fig 4: Diffraction arm\end{center}

The lamp is initially centered, though you can adjust it if needed. Make sure
the filament is vertical for a finer spectral line on the diffraction grating.

You can set a convergence angle of the beam by moving the lens. Untighten a
wing nut (WN), slide the lens holder and tighten it again (do not overtighten
it, a slight torque is enough to prevent the lens from an accidental movement).

\textbf{Powering the lamp and lasers}

To power light sources plug in the DC adapter into the outlet on the table.
Connect the adapter to the socket of the lamp or the laser. Polarity is
important for lasers and it is set properly initially.

\textbf{Lasers}

You are provided with three lasers: red ($\lambda = 659$\,nm), green
($\lambda = 530$\,nm) and blue ($\lambda = 400$\,nm), marked with color tape.

\noindent\fbox{\parbox{\dimexpr\linewidth-2\fboxsep-2\fboxrule\relax}{\centering\large
Lasers are dangerous! Do not point lasers at your eyes or look into the laser
beam under any circumstances! Beware the reflected beams!}}

Red and blue lasers share the powering attachment. To switch from red to blue
laser unscrew the attachment from the red laser and screw it to the blue one
(fig. 5a).

\begin{center}\includegraphics[width=0.9\linewidth]{figures/APhO_2017_Q4_fig5.jpg}\end{center}

\begin{center}Fig 5: (a) Switching the powering attachment. (b) Adjusting the laser beam.\end{center}

In photometric experiments put the button lock on the laser, place the laser
onto the holder and fix it with the rubber band (fig. 5b). Put the holder disk
into the left hole of the main bench. Rotate the button lock to turn on the
laser (1).

You may need to adjust the laser beam. Start with rotating the laser around its
axis until the beam lies in the horizontal plane (2). Then rotate the holder
disk to align the beam with the optic bench (3).

\textbf{Photodiode}

For photometric experiments you will need lasers and photodiode to measure the
transmitted intensity of the light. The photodiode holder has legs that may be
inserted in the slots of the main bench (fig. 6a).

\begin{center}\includegraphics[width=0.9\linewidth]{figures/APhO_2017_Q4_fig6.jpg}\end{center}

\begin{center}Fig. 6: (a) Installation of the photodiode. (b) Connection of the multimeter.\end{center}

Photodiode is attached to holder with magnets and can be adjusted for the laser
beam.

To conduct measurements with the photodiode connect it to the multimeter in DC
current mode (μA) (fig. 6b). The current of the photodiode is proportional to
the intensity of the incident light, so we will measure the light intensity in
μA.

\textbf{Photonic crystal samples}

\noindent\fbox{\parbox{\dimexpr\linewidth-2\fboxsep-2\fboxrule\relax}{\centering\large
Photonic crystals are very fragile! Do not touch the crystals! Hold the samples
only by the frame!}}

You are provided with three samples of different photonic crystals. They have
frontfaces of different colors and a letter, sample name, (X, Y or Z) on the
top. Make sure that when you place the sample into the holder, the light beam
falls onto the frontface and the letter is on top. While conducting
measurements, place the holder disk into the right hole of the main bench. Make
sure to rotate the holder disk clockwise, so the reflected light beam is
averted.

\begin{center}\includegraphics[width=0.6\linewidth]{figures/APhO_2017_Q4_fig7.jpg}\end{center}

\begin{center}Fig 8: Samples of photonic crystals\end{center}

\textbf{Bragg-Snell law}

\begin{center}\includegraphics[width=0.81\linewidth]{figures/APhO_2017_Q4_fig8.png}\end{center}

\begin{center}Fig. 8. (a) Photonic crystall structure. (b) Refractive index $n(z)$ dependence.\end{center}

Photonic crystal samples being explored in this problem consist of layers with
different refractive indices $n_i$. The refractive index changes periodically
along $z$ axis with period $D$ and does not depend on a wavelength $\lambda$.
Let us denote the overall mean refractive index $n$ and refractive indices
difference $\Delta n = n_{\text{max}} - n_{\text{min}}$. In AAO crystals
\begin{equation}
\Delta n \ll n.
\tag{1}
\end{equation}

We consider parallel monochromatic incident beam with the wavelength $\lambda$
and constant intensity falling on the photonic crystal. The incidence angle is
$\theta$. The beams reflected from different layers interfere. As a result,
mirror reflection has maximums at angles $\theta$ that satisfy the equation:
\begin{equation}
2D\sqrt{n^2 - \sin^2\theta} = m\lambda,
\tag{2}
\end{equation}
where $m = 1, 2, \ldots$ is an integer number representing the interference
order.

Equation (2) is called the \textbf{Bragg-Snell law}. Transmittance minimums are
observed at the same angles $\theta_i$. If angle $\theta$ is constant, and
wavelength $\lambda$ varies, Bragg-Snell law may be regarded as an equation for
$\lambda$. Fig.9 shows sketches of reflectance and transmittance spectra. Each
transmittance minimum corresponds to some integer $m$ value according to the
Bragg-Snell law.

\begin{center}\includegraphics[width=0.61\linewidth]{figures/APhO_2017_Q4_fig9.png}\end{center}

\begin{center}Fig. 9: Reflectance and transmittance spectra.\end{center}

The wavelengths satisfying the Bragg-Snell law (2) will be called the
\textbf{wavelengths of transmittance minimums}. The wavelengths of
transmittance minimums depend on the incidence angle $\theta$. The
transmittance minimums wavelengths at $\theta = 0$ will be called the
\textbf{normal wavelengths of transmittance minimums} for a particular photonic
crystal.

\noindent\fbox{\parbox{\dimexpr\linewidth-2\fboxsep-2\fboxrule\relax}{\centering\large
You do not need to evaluate errors in this problem.}}

\textbf{Part A. Sample X. Spectral measurements (3.5 points)}

Sample X has a simple structure (fig. 10). Its period consists of two layers
with the same thickness $D_X/2$ and with close refractive indices $n_1$ and
$n_2$:
\begin{equation}
n_1 - n_2 = \Delta n \ll n_X = \frac{n_1 + n_2}{2}
\tag{3}
\end{equation}

\begin{center}\includegraphics[width=0.25\linewidth]{figures/APhO_2017_Q4_fig10.png}\end{center}

\begin{center}Fig. 10: The structure of Sample X .\end{center}

\begin{center}\includegraphics[width=0.91\linewidth]{figures/APhO_2017_Q4_fig11.png}\end{center}

Fig 11. Experimental setup for spectral measurements: (F) lens, (L) lamp,
located in the focal point of the lens, (DG) diffraction grating, (PhCr)
photonic crystal.

1. Assemble the \textbf{experimental setup for the spectral measurements}.
Place the diffraction arm into the left hole of the main bench;

2. Turn on the lamp. Look at the diffraction grating from the other end of the
optical bench. Rotate the diffraction arm until you see the first diffraction
order (a rainbow strip).

3. Install the sample holder into the right hole of the main bench (set the
triangular mark on zero). Place Sample X in the holder (the light should fall
onto the red frontface). Look at the rainbow strip of the diffraction grating
through the sample.

4. If you rotate the sample clockwise at the right angle, you can see a dark
region in the transmitted spectrum (fig. 12). This is the \textbf{transmittance
minimum} corresponding to $m = 1$ in the Bragg-Snell law. When you rotate the
sample, the wavelength of the transmittance minimum changes.

\begin{center}\includegraphics[width=0.41\linewidth]{figures/APhO_2017_Q4_fig12.jpg}\end{center}

Fig. 12. Photo of the transmittance spectrum of the sample X. White arrow shows
the transmittance minimum (a dark area in a rainbow strip).

5. By moving the lens, you can see wider or narrower part of the spectrum. For
precise measurements place the lens in such a manner, that the lamp will be
located in the focal point.

\textbf{A.1}\quad Light falls on the diffraction grating perpendicular to its
surface. Write down the equation, connecting the diffraction angle $\varphi$
and the wavelength $\lambda$. Grating period is $h = 1000$\,nm. You can write
down the answer without deriving it. \hfill 0.1pt

\textbf{A.2}\quad Measure the angle of incidence $\theta$, at which you see the
transmittance minimum, as a function of the diffraction angle $\varphi$. Take
reading at such $\theta$ that transmittance minimum (dark region) can be seen
in the middle of the grating window. \hfill 1.0pt

\textbf{A.3}\quad Find new coordinates instead of $\theta$ on $\varphi$ to test
that graph would be linear. Plot the graph in new coordinates. \hfill 1.5pt

\textbf{A.4}\quad Determine the sample period $D_X$ and the mean refractive
index $n_X$. \hfill 0.9pt

\textbf{Part B. Sample X. Laser measurements (5.0 points)}

The porosity of the crystal $p$ is a fraction of the crystal volume filled with
air (channels). Crystal effective refractive index can be expressed in terms of
air refractive index $n_a = 1$ and aluminum oxide refractive index $n_{AAO}$:
\begin{equation}
n_{dry} = \sqrt{p n_a^2 + (1 - p) n_{AAO}^2}.
\tag{4}
\end{equation}

If channels are filled with water, the refractive index of crystal changes:
\begin{equation}
n_{wet} = \sqrt{p n_w^2 + (1 - p) n_{AAO}^2},
\tag{5}
\end{equation}
where $n_w = 1.33$ is the refractive index of water.

In this part we will determine the porosity $p$ of the photonic crystal by
comparing the experimental data for dry and wet samples. We will use the setup
for laser measurements.

\textbf{B.1}\quad Choose the laser such that the transmittance minimum may be
observed at laser wavelength $\lambda$ at some sample rotation angle $\theta$.
Write down the wavelength of the laser you have chosen. \hfill 0.1pt

Use the chosen laser in all tasks of part B.

Assemble the \textbf{experimental setup for laser measurements} (fig. 13).

\noindent\fbox{\parbox{\dimexpr\linewidth-2\fboxsep-2\fboxrule\relax}{\centering\large
Do not stare into the laser beam directly! Keep track of reflected beams, they
must not get into your eyes. Even short exposure may seriously damage your
eyes. Do not hold your head on beam level when you conduct any laser
experiments.}}

Before collecting any data in laser experiment, make sure that laser intensity
is constant. If it is not, wait for 3-5 minutes.

\begin{center}\includegraphics[width=0.91\linewidth]{figures/APhO_2017_Q4_fig13.png}\end{center}

Fig. 13. Setup for laser measurements: (Laser) laser, (PhCr) photonic crystal,
(PhD) photodiode, connected to an amperemeter. The current measured by the
amperemeter is proportional to the intensity of light falling on the
photodiode.

\textbf{B.2}\quad Measure the intensity of laser light $I_t$ (in μA)
transmitted through the photonic crystal as a function of $\theta$.
\hfill 1.0pt

\textbf{B.3}\quad Plot the graph $I_t(\theta)$. \hfill 1.0pt

\textbf{B.4}\quad Determine the angle of incidence $\theta_1$, corresponding to
the transmittance minimum at the laser wavelength. Determine the width of the
valley $\Delta\theta_1$ on plot $I_t(\theta)$ at the level of half-depth (fig.
14a). \hfill 0.2pt

\begin{center}\includegraphics[width=0.72\linewidth]{figures/APhO_2017_Q4_fig14.png}\end{center}

Fig. 14. Width of the transmittance valley at half-depth for the dependence on
(a) incident angle, (b) on wavelength.

\textbf{B.5}\quad Using the Bragg-Snell law (2), determine the transmittance
minimum normal wavelength $\lambda_X$ for sample X. Use the value $n_X$ from
part A, $\theta_1$ from B.4 and laser wavelength $\lambda$. \hfill 0.2pt

When light falls perpendicular to the crystal ($\theta = 0$), spectral width
$\Delta\lambda$ (fig. 14b) of transmittance minimum, corresponding to $m = 1$,
can be expressed in terms of $n$ and $\Delta n$ as follows:
\begin{equation}
\frac{\Delta\lambda}{\lambda} = \frac{2}{\pi}\frac{\Delta n}{n}.
\tag{6}
\end{equation}

To estimate $\Delta n_X$ for the sample X, assume that relational spectral
width of transmittance minimum $\frac{\Delta\lambda}{\lambda}$ does not change
with rotation.

\textbf{B.6}\quad Estimate refractive indices difference $\Delta n_X$ of sample
X. \hfill 0.6pt

Place some water on the face side of the sample X. Carefully place a coverslip
over the sample to prevent water evaporation during the experiment. There is a
solid oxide layer 15\,nm thick on the back side of the sample which prevents
water from pouring out.

\noindent\fbox{\parbox{\dimexpr\linewidth-2\fboxsep-2\fboxrule\relax}{\centering\large
Be careful! Samples are very fragile. Their width is a quarter of a hair
diameter. Operate samples only by the frame. The coverslip is very fragile too!
Hold it only with tweezers.}}

\textbf{B.7}\quad Determine the angle of incidence $\theta_2$ corresponding to
the transmittance minimum of the wet sample at the laser wavelength.
\hfill 0.3pt

\textbf{B.8}\quad Determine the mean porosity $p_X$ of the sample X and
aluminum oxide refractive index $n_{AOO}$. \hfill 1.0pt

\textbf{B.9}\quad Determine the mean porosities $p_1$ and $p_2$ of the layers
in sample X. \hfill 0.6pt

\textbf{Part C. Sample Y. Several transmittance minimums (4.5 points)}

Sample Y structure is more complicated than the one of sample X. Sample Y has 4
transmittance minimums in the visible spectral range (400--800\,nm)
corresponding to 4 consecutive integers $m$ in the Bragg-Snell law (2). The
refractive indices of samples Y and X are the same $n_Y = n_X$.

\textbf{C.1}\quad Although there are 4 transmittance minimums for sample Y, you
can observe only three of them in spectral measurements. Using spectral
measurements (see part A), determine the normal wavelengths of these three
transmittance minimums $\lambda_1^{sp}$, $\lambda_2^{sp}$, $\lambda_3^{sp}$.
\hfill 0.6pt

Now conduct laser measurements for the sample Y (see part B). Laser
measurements let you determine normal wavelengths of transmittance minimums
with greater accuracy. Moreover, you will be able to determine transmittance
values which you will need in part E.

\textbf{C.2}\quad For red laser, measure transmitted light intensity $I_{red}$
as a function of the sample rotation angle $\theta$. \hfill 0.5pt

\textbf{C.3}\quad For green laser, measure transmitted light intensity
$I_{green}$ as a function of the sample rotation angle $\theta$. \hfill 0.5pt

\textbf{C.4}\quad For blue laser, measure transmitted light intensity
$I_{blue}$ as a function of the sample rotation angle $\theta$. \hfill 0.5pt

\textbf{C.5}\quad Using data obtained in tasks C.2--C.4, determine normal
wavelengths of 4 transmittance minimums. \hfill 0.6pt

\textbf{C.6}\quad Determine integers $m$ corresponding to 4 normal wavelengths
obtained in C.5. You may plot a graph if you need. \hfill 1.0pt

\textbf{C.7}\quad Determine the period $D_Y$ of the sample Y in nanometers.
\hfill 0.2pt

\begin{center}\includegraphics[width=0.52\linewidth]{figures/APhO_2017_Q4_fig15.png}\end{center}

\begin{center}Fig. 15. Determining transmittance $t$ with use of $I(\theta)$ plot\end{center}

\textbf{C.8}\quad Determine the transmittance values $t$ of 4 transmittance
minimums of the sample Y (fig. 15). \hfill 0.6pt\\
\textit{These values will be used in part E to determine the structure of the
sample Y. You may plot graphs if you need it, though it is not required.}

\textbf{Part D. Sample Z. Missed transmittance minimums (4.5 points)}

The Bragg-Snell law (2) determines angles and wavelengths at which
transmittance minimum may be observed. Transmittance values $t$ depend on the
internal structure of the period. A sample may have a specific structure such
that some transmittance minimums predicted by the Bragg-Snell law have $t = 1$
and therefore can not be seen. We will call these transmittance minimums
\textbf{missed transmittance minimums}. We will call minimums that are not
missed \textbf{visible}.

Regard transmittance minimum as missed if it can not be seen in spectral or
laser experiments at the position predicted by the Bragg-Snell law. The sample
Y has no missed transmittance minimums. Conversely, the sample Z has 2 missed
minimums in the visible spectrum. This means that not all visible transmittance
minimums correspond to consecutive $m$.

Mean refractive indices of the samples X and Z are the same: $n_Z = n_X$.

\textbf{D.1}\quad Determine normal wavelengths $\lambda_Z$ of visible
transmittance minimums of the sample Z. Use spectral and laser experiments.
Describe your experiments with sketches and equations. \hfill 1.2pt

\textbf{D.2}\quad Determine integers $m$ corresponding to visible transmittance
minimums. You may plot a graph if you need. \hfill 2.0pt

\textbf{D.3}\quad Determine period $D_Z$ of the sample Z in nanometers.
\hfill 0.3pt

\textbf{D.4}\quad Determine, which transmittance minimums are missing. Derive
their $m'$ in Bragg-Snell law and normal wavelengths $\lambda'_Z$.
\hfill 1.0pt

\textbf{Part E. Samples Y and Z. Internal structure of the period. (2.5 points)}

The Bragg-Snell law (2) describes angles and wavelengths, where reflectance
maximums and transmittance minimums may be observed (fig. 9). Transmittance
values depend on the internal structure of the period.

Simple theory for the reflectance values is discussed in this part. Normal
incidence of light ($\theta = 0$) is considered. Such being the case, the
Bragg-Snell law simplifies:
\begin{equation}
2Dn = m\lambda
\tag{7}
\end{equation}
where $D$ is the period of the crystal.

In the model we assume that light reflects only from boundaries where it
travels from high to low refractive index layer. Reflective boundaries
considered in the model are marked in fig. 16 with arrows.

The samples X, Y, and Z consist of the layers with the same thickness and 2
different refractive indices, so we will limit ourselves only to this
particular case.

Let $d_l$ be the thickness of bilayer, $D$ crystal period, $\delta_j$ the
distance between the beginning of the period and the reflective boundary $j$.

\begin{center}\includegraphics[width=0.41\linewidth]{figures/APhO_2017_Q4_fig16.png}\end{center}

Fig.16. Example structure of repeating crystal period $D$. Dark green layers
have high refractive index, bright yellow layers have low refractive index. All
layers are of the same thickness $d_l/2$. The reflective boundaries considered
in the model are marked with arrows.

For example, structure on Fig. 16 has crystal period $D = 3.5d_l$, the
distances to reflective boundaries are $\delta_1 = 0$, $\delta_2 = d_l$,
$\delta_3 = 2.5d_l$.

For the beam reflected from layer $j$ phase advance is
\begin{equation}
\varphi_{j_m} = 2\pi\frac{2\delta_j n}{\lambda} = 2\pi m\frac{\delta_j}{D}
\tag{8}
\end{equation}

Considering interference between beams reflected from one period, we deduce
that reflective intensity for reflectivity maximum $m$ of the Bragg-Snell law
is:
\begin{equation}
I_{refl,m} \sim \left|\sum_j \exp(i\varphi_{j_m})\right|^2
\tag{9}
\end{equation}
where $i$ -- imaginary unit.

The model considered above is not quantitatively correct. But it keeps the
order of reflectance maximums: you can compare the magnitudes of maximums and
say, which would be more intensive. If reflectance maximum is intense,
transmittance value $t$ of transmittance minimum is low. If reflectance maximum
is low enough, we consider it as missed transmittance minimum.

In the appendix there is a table with the values $I_m/I_0$ calculated according
to equation (9) for the integer $m$ from 1 to 20. The intensity $I_0$ is
calculated according to (9) for $m = 0$. In that case all the reflected beams
interfere without phase shift to output the maximum intensity possible. Values
$I_m/I_0$ are expressed in $\%$.

The structures of the samples Y and Z are in the table among others.

\textbf{E.1}\quad Compare transmittance values for the sample Y obtained in C.8
with the table in the appendix. Determine the structure of the sample Y. Write
the name of the structure in the answer sheet. \hfill 1.2pt

\textbf{E.2}\quad Determine the structure of the sample Z. Write the name of
the structure in your answer sheet. \hfill 1.3pt

\textbf{Appendix}

Possible variants of structures of samples Y and Z. Names of the structures are
given under the pictures. Numbers in the table are $I_m/I_0$.

% The table is printed rotated by 90 degrees (landscape) in the source, as a raster image; it is set upright here.
\begin{center}
\setlength{\tabcolsep}{3pt}
\begin{tabular}{c*{8}{c}}
 & & & & & & & \multicolumn{2}{c}{$d_1/d_2 = 1.13$} \\
 & \includegraphics[angle=90,scale=0.8]{figures/APhO_2017_Q4_fig17.png}
 & \includegraphics[angle=90,scale=0.8]{figures/APhO_2017_Q4_fig18.png}
 & \includegraphics[angle=90,scale=0.8]{figures/APhO_2017_Q4_fig19.png}
 & \includegraphics[angle=90,scale=0.8]{figures/APhO_2017_Q4_fig20.png}
 & \includegraphics[angle=90,scale=0.8]{figures/APhO_2017_Q4_fig21.png}
 & \includegraphics[angle=90,scale=0.8]{figures/APhO_2017_Q4_fig22.png}
 & \includegraphics[angle=90,scale=0.8]{figures/APhO_2017_Q4_fig23.png}
 & \includegraphics[angle=90,scale=0.8]{figures/APhO_2017_Q4_fig24.png} \\
\midrule
\textbf{m} & \textbf{i5-5} & \textbf{i3-3} & \textbf{i2-1} & \textbf{n-3} & \textbf{n-6} & \textbf{n-9} & \textbf{hi5-5} & \textbf{h5-5} \\
\midrule
1 & 1 & 3 & 11 & 3 & 1 & 0 & 1 & 0 \\
2 & 0 & 0 & 33 & 7 & 1 & 0 & 0 & 0 \\
3 & 1 & 6 & 11 & 56 & 1 & 0 & 1 & 0 \\
4 & 0 & 0 & 33 & 56 & 2 & 0 & 0 & 0 \\
5 & 2 & 41 & 11 & 7 & 6 & 1 & 2 & 0 \\
6 & 0 & 0 & 100 & 3 & 48 & 1 & 1 & 0 \\
7 & 5 & 41 & 11 & 100 & 48 & 2 & 3 & 0 \\
8 & 0 & 0 & 33 & 3 & 6 & 5 & 6 & 0 \\
9 & 41 & 6 & 11 & 7 & 2 & 45 & 36 & 11 \\
10 & 0 & 0 & 33 & 56 & 1 & 45 & 1 & 72 \\
11 & 41 & 3 & 11 & 56 & 1 & 5 & 32 & 15 \\
12 & 0 & 100 & 100 & 7 & 1 & 2 & 10 & 2 \\
13 & 5 & 3 & 11 & 3 & 100 & 1 & 1 & 0 \\
14 & 0 & 0 & 33 & 100 & 1 & 1 & 2 & 0 \\
15 & 2 & 6 & 11 & 3 & 1 & 0 & 0 & 0 \\
16 & 0 & 0 & 33 & 7 & 1 & 0 & 1 & 0 \\
17 & 1 & 41 & 11 & 56 & 2 & 0 & 0 & 0 \\
18 & 0 & 0 & 100 & 56 & 6 & 0 & 6 & 7 \\
19 & 1 & 41 & 11 & 7 & 48 & 100 & 28 & 27 \\
20 & 100 & 0 & 33 & 3 & 48 & 0 & 25 & 23 \\
\bottomrule
\end{tabular}
\end{center}
